%!TEX program = lualatex
\documentclass[ngerman,headheight=70pt]{scrartcl}
\usepackage[ngerman]{babel}
\usepackage{blindtext}
\usepackage[utf8]{luainputenc}
\usepackage{scrlayer-scrpage}
\usepackage[T1]{fontenc}
\usepackage{fontspec}
\usepackage{multicol}
\usepackage{enumitem}
\usepackage{epigraph}
\usepackage[german=quotes]{csquotes}
\usepackage[bottom=2cm,footskip=8mm]{geometry}
\setmainfont{Source Sans Pro}
\MakeOuterQuote{"}

\setlength{\parindent}{0pt}
\setlength{\parskip}{4pt}

\pagestyle{plain.scrheadings}
\thispagestyle{scrheadings}
\lohead{\vspace*{-1.5cm}\includegraphics[scale=0.4]{../images/stupa_logo_rot.pdf}}
\rohead{\\Präsidium des\\Studierendenparlaments}

\lofoot[Universität Hamburg · Präsidium des Studierendenparlaments\\\hrule Von-Melle-Park 5 · D-20146 Hamburg]{Universität Hamburg · Präsidium des Studierendenparlaments\\\hrule Von-Melle-Park 5 · D-20146 Hamburg}
\cofoot[]{}
\cohead[\thepage]{}

\renewcommand{\thesubsection}{TOP \arabic{subsection}}
\renewcommand{\thesubsubsection}{\arabic{subsubsection}.}
\setcounter{subsection}{-1}

\begin{document}
    UHH · StuPa-Präsidium · Von-Melle-Park 5 · D-20146 Hamburg

    \section*{Beschluss des Studierendenparlamentes vom 1. Dezember 2016}
    \subsection*{Für einen fairen Wahlkampf}

    Das Studierendenparlament ruft alle Kandidierenden aller Listen dazu auf,
    den Wahlkampf zur Wahl des Studierendenparlaments 2017/2018 auf eine
    sachlich-argumentative Weise und mit gegenseitigem Respekt durchzuführen und
    auf platte Polemik, Beleidigungen und Verleumdungen zu verzichten.

    Alle Listen sind dazu aufgerufen, ihr Bestmöglichstes zu tun um die
    Wähler*innen über ihre Ziele und Inhalte zu informieren. Auch ist auf den
    Wahrheitsgehalt von Wahlversprechen zu achten.

    Die Listen werden dazu angehalten, von jeglicher Ausgabe von Waffeln, Glühwein
    oder anderer Geschenke im direkten Austausch gegen die Stimmabgabe für eine
    spezifische Liste abzusehen und auch andere Arten von Stimmkäufen zu
    unterlassen.

    Die Listen werden dazu aufgerufen auf ein neutrales Umfeld an den Wahlurnen
    zu achten.

    Dies beinhaltet das Absehen von Listenwerbung im Umfeld von 5m um die
    Wahlurnen, das Hinterlegen von Informationsmaterial oder das Tragen von
    listenzugehörigen Kleidungsstücken an der Urne, wie auch das direkte an der
    Urne "Abliefern" von potentiellen Wähler*innen. Die Wahlhelfer*innen des
    Präsidiums verpflichten sich, die Einhaltung dessen überparteilich zu wahren.

    Des Weiteren sieht das Studierendenparlament alle Kandidierenden in der
    Verantwortung, jeglichen Versuch und Vollendung von Wahlfälschung
    zurückzuweisen. Alle Listen werden dazu aufgefordert, bei der Wahl und der
    Auszählung auf die Einhaltung demokratischer und rechtlicher Grundlagen zu
    achten.

    Für einen zivilisierten, fairen und inhaltlichen Wahlkampf spricht sich das
    Studierendenparlament für folgende Punkte aus:

    \begin{itemize}
        \item Die Auseinandersetzungen zwischen den verschiedenen Listen während
            der StuPa-Wahl sind politisch-inhaltlich und argumentativ im Streit um
            studentische Interessenvertretung und ihren Bezug zu gesellschaftlichen
            Auseinandersetzungen zu führen.
        \item Die Kritik an den Positionen anderer muss authentisch, sachlich und
            überprüfbar sein. Polemik und ironische Mittel beschränken sich auf
            den exemplarischen Gehalt vertretener politischer Positionen.
        \item Wahlhelfer*innen und Präsidium nehmen im Rahmen ihrer Tätigkeit
            ausschließlich allgemeine, nicht "listenbornierte" Aufgaben der VS
            wahr. Davon unbenommen sind Wahlkampfaktivitäten neben der
            "Amtstätigkeit".
        \item Es wird darauf geachtet, dass alle Listen durch andere nicht in
            ihrer (Öffentlichkeits-)Arbeit eingeschränkt werden. Die Wahlkämpfenden
            achten darauf, dass auch die Materialien der anderen Listen pfleglich
            behandelt werden.
        \item Das Studierendenparlament führt zur erweiterten Mobilisierung zwei
            Veranstaltung zur Listenvorstellung durch. Alle Listen beteiligen sich
            an der Ankündigung und Mobilisierung für die Veranstaltung, um einen
            möglichst hohen Beteiligungsgrad zu erzielen.
        \item Alle Listen werden gebeten, diesen Beschluss bei Infotischen und
            Veranstaltungen mit auszulegen.
    \end{itemize}

\end{document}
